Die Temperaturen der Zustände folgen aus $T/T_0 = (p/p_0)\cdot(V/V_0)$:
\begin{center}
\begin{tabular}{lcccc}
 & A & B & C & D \\ \hline
 $p$ in $p_0$ & 4 & 4 & 2 & 1 \\
 $V$ in $V_0$ & 1 & 2 & 4 & 4 \\
 $T$ in $T_0$ & 4 & 8 & 8 & 4
\end{tabular}
\end{center}
\begin{abcliste}
\abc A $\to$ B ist $\result{\text{isobar}}$ (waagrecht, Expansion), B $\to$ C $\result{\text{isotherm}}$ (Hyperbel, $T = 8\,T_0$, Expansion), C $\to$ D $\result{\text{isochor}}$ (senkrecht, Abkühlung), D $\to$ A $\result{\text{isotherm}}$ (Hyperbel, $T = 4\,T_0$, Kompression).

\abc
\begin{enumerate}
\item $\result{\text{Richtig}}$: Auf einer Isotherme ist $pV$ konstant; $p$ fällt von $4\,p_0$ auf $2\,p_0$, $V$ steigt von $2\,V_0$ auf $4\,V_0$.
\item $\result{\text{Falsch}}$: $T_C = 8\,T_0$, $T_A = 4\,T_0$. C liegt auf der äusseren Hyperbel.
\item $\result{\text{Richtig}}$: Isobar ändern sich $V$ und $T$ um denselben Faktor; $V$ verdoppelt sich, also auch $T$ (von $4\,T_0$ auf $8\,T_0$).
\item $\result{\text{Falsch}}$: Je weiter eine Isotherme vom Ursprung entfernt ist, desto grösser ist $pV$ und damit $T$. Die Hyperbel durch D und A gehört zu $4\,T_0$, jene durch B und C zu $8\,T_0$.
\item $\result{\text{Richtig}}$: Isochor ändern sich $p$ und $T$ um denselben Faktor; $p$ halbiert sich, also auch $T$ (von $8\,T_0$ auf $4\,T_0$).
\item $\result{\text{Falsch}}$: Das Gas wird komprimiert, aber isotherm: Der Druck steigt dabei von $p_0$ auf $4\,p_0$.
\end{enumerate}
\end{abcliste}
